\section{Welfare}\label{sec:welfare}

The welfare cost of the set-aside is the structural object that disciplines the policy comparison. Under the main specification (Assumption~\ref{a:setaside}, with $\hat F_c^{\text{SME,Post}}$ estimated from post-period exits), that cost is \welfLossPctLthirtyNp~percent of $p_{S_1}$ in non-pharmaceuticals and \welfLossPctLthirtyPh~percent in pharmaceuticals at $\lambda = \lambdaMain$, where $\lambda$ is the marginal cost of public funds (formal definition below). The price effect is the beginning of that account, not its end. Allocative waste from routing the contract to a higher-cost SME winner, tax-distortion loss on the additional government outlay, and the implicit transfer to SME bidders combine in a single-equation decomposition (eq.~\ref{eq:welfare}) that the structural framework identifies separately. The same structure that delivered the decomposition in Section~\ref{sec:results} explains why the welfare comparison is stable in one class and model-sensitive in the other.

The price effect is one piece of that welfare cost, and one of three. Two others matter and are routinely left out of reduced-form accountings. \emph{Allocative waste} arises because the set-aside routes the contract to an SME winner whose cost may exceed the lowest cost in the full pool; the extra production cost is lost to society, not a transfer. \emph{Tax-distortion loss} arises because the additional outlay that the government makes is financed from distortionary taxation, so each marginal real costs $(1+\lambda)$ in welfare shadow terms. The transfer from government to SMEs---the government's extra outlay $\Delta_{\text{gov}} = p_{S_3} - p_{S_1}$ minus the allocative wedge---is zero-sum in a social-planner ledger and is not destruction. This section isolates the two destruction pieces.

Formally, per-auction welfare loss of the SME-only set-aside relative to the open regime is
\begin{equation}\label{eq:welfare}
\mathrm{Loss} \;=\; \underbrace{c_{(1)}^{S_3} - c_{(1)}^{S_1}}_{\mathrm{DWL}_{\text{alloc}}}
\;+\; \underbrace{\lambda \cdot (p_{S_3} - p_{S_1})}_{\mathrm{MCPF\ distortion}}.
\end{equation}
The first term is the allocative wedge. It measures by how much the winner's production cost is higher in the restricted regime than in the open regime. The second term is the tax-distortion loss on the government's extra outlay $\Delta_{\text{gov}} = p_{S_3} - p_{S_1}$, scaled by the marginal cost of public funds $\lambda$. The natural normalization is $\mathrm{Loss} / p_{S_1}$, and I report it that way.

\subsection{Point estimates}

Table~\ref{tab:v3_welfare} gives the decomposition at $\lambda = \lambdaMain$, the \citet{ballard1985} benchmark; \citet{hendren2020} report a comparable range on the modern MVPF basis across U.S.\ tax-financed programs, and \citet{finkelstein2020} discuss how MVPF connects to the marginal-cost-of-public-funds framework used here. The reported MCPF sensitivity in this paper covers the relevant interval. In non-pharmaceuticals, $\Delta_{\text{gov}} = \welfDeltaGovNp$ of the reference price, of which $\mathrm{DWL}_{\text{alloc}} = \welfDwlAllocNp$ is allocative waste and the remaining $\welfTransferNp$ is an implicit transfer to SME bidders. MCPF distortion adds $\lambdaMain \times \welfDeltaGovNp \approx \welfMcpfDistLthirtyNp$. Total welfare loss is $\mathrm{DWL}_{\text{alloc}} + \mathrm{MCPF\ dist.} = \welfDwlAllocNp + \welfMcpfDistLthirtyNp = \welfTotalLossLthirtyNp$, or \welfLossPctLthirtyNp~percent of $p_{S_1}$.

In pharmaceuticals, the same accounting is larger on every margin. $\Delta_{\text{gov}} = \welfDeltaGovPh$; $\mathrm{DWL}_{\text{alloc}} = \welfDwlAllocPh$; implicit transfer $\welfTransferPh$; MCPF distortion $\lambdaMain \times \welfDeltaGovPh \approx \welfMcpfDistLthirtyPh$; total welfare loss $\approx \welfTotalLossLthirtyPh$, or \welfLossPctLthirtyPh~percent of $p_{S_1}$. The allocative component dominates in both classes, but especially in pharma, where a thin SME pool and higher goods heterogeneity make bidder exclusion more likely to route the contract to a high-cost protected winner.\footnote{Algebraic identity: $\mathrm{Loss}/p_{S_1} = (\mathrm{DWL}_{\text{alloc}} + \lambda \cdot \Delta_{\text{gov}})/p_{S_1} = (1 + \lambda) \cdot \Delta_{\text{gov}}/p_{S_1} - \mathrm{transfer}/p_{S_1}$. The pharma welfare cost (\welfLossPctLthirtyPh~percent of $p_{S_1}$) is close to the structural price ratio $\Delta_{\text{gov}}/p_{S_1} \approx \welfPriceRatioPharm$ percent because the implicit SME-winner transfer ($\welfTransferPh$ of $p_{S_1}$) happens to be near $\lambda \cdot \Delta_{\text{gov}} = \lambdaMain \times \welfDeltaGovPh \approx \welfMcpfDistLthirtyPh$, leaving residual $\approx (1)\Delta_{\text{gov}}/p_{S_1}$. In non-pharma the transfer ($\welfTransferNp$) exceeds $\lambda \cdot \Delta_{\text{gov}} = \welfMcpfDistLthirtyNp$, so the welfare cost (\welfLossPctLthirtyNp~percent) lies materially below the price ratio (\welfPriceRatioNpFn~percent); the two objects should not be conflated.}

\subsection{Endogenous entry as a partial offset in welfare units}

The fixed-pool-vs-realized attenuation documented in Section~\ref{sec:results} (\latentShockRatioNp~percent in non-pharmaceuticals, \latentShockRatioPh~percent in pharmaceuticals) carries over to the welfare formula: both terms in equation~\eqref{eq:welfare} load on the simulated price margin, so a fixed-pool welfare calculation would overstate the realized burden by the same factor. The participation response is therefore part of the welfare object, not separable from it---a fixed-pool counterfactual misstates the policy cost by exactly the latent-shock attenuation, and the V3-vs-V0 ranking in non-pharmaceuticals is preserved either way.

\subsection{Confidence intervals}

Cluster bootstrap with $\bootReps$ replicates (Table~\ref{tab:v3_welfare_ci}) delivers 95 percent CIs on $\mathrm{Loss}/p_{S_1}$ at $\lambda = \lambdaMain$ summarized in Table~\ref{tab:welfare_ci_summary}.

\begin{table}[!htbp]
\centering
\begin{threeparttable}
\caption{Cluster-bootstrap CIs on welfare loss at $\lambda = 0.30$.}
\label{tab:welfare_ci_summary}
\small
\setlength{\tabcolsep}{6pt}
\begin{tabular}{lcc}
\toprule
Class      & Bootstrap mean  & 95\% CI \\
\midrule
Non-pharma & \welfLossPctLthirtyNp\% & $[\welfLossPctLoNp\%,\ \welfLossPctHiNp\%]$ \\
Pharma     & \welfLossPctLthirtyPh\% & $[\welfLossPctLoPh\%,\ \welfLossPctHiPh\%]$ \\
\bottomrule
\end{tabular}
\begin{tablenotes}\footnotesize
\item Cluster bootstrap at the auction level, $B = 500$ replicates, all-bidders UH-clean $F_c$. Welfare loss reported as percent of the open-regime baseline $p_{S_1}$ at $\lambda = \lambdaMain$. Full $\lambda$-grid CIs in Table~\ref{tab:v3_welfare_ci}.
\end{tablenotes}
\end{threeparttable}
\end{table}

The intervals leave no doubt about the ordering. Even the lower bound in pharma is above the non-pharma point estimate. The class heterogeneity is therefore not sampling noise around a common welfare effect; it is a structural feature of the protected pool and the good being procured. Figure~\ref{fig:v3_welfare_forest} reports the CIs across the $\lambda$ grid and overlays the strict-invariance benchmark.

\begin{figure}[!htb]
\centering
\includegraphics[width=0.98\textwidth]{../output/figures/fig_v3_welfare_forest.pdf}
\caption[Welfare loss intervals by class and MCPF $\lambda$.]{\textit{Welfare loss as \% of $p_{S_1}$, by class and MCPF $\lambda$.} Circles and horizontal whiskers report the cluster-bootstrap mean and 95\% CI under the main specification ($\bootReps$ replicates). Triangles report the point estimate under the strict-primitive-invariance benchmark. The dotted vertical line marks the \didBenchmarkPct\% benchmark implied by the difference-in-differences coefficient of \ref{app:did}: the \structuralWindowM-month log-price effect of $\didLowerDec$, exponentiated and adjusted for the average pre-period price-to-reference ratio of about \didPreRefRatio, gives a normalized benchmark of approximately \didBenchmarkPct\% of $p_{S_1}$, against which the structural welfare cost is reported. Both specifications place the structural welfare cost two to four times the DiD benchmark, and neither CI crosses the benchmark estimate. In pharma, the strict benchmark lies below the main specification because it attributes less of $\Delta_{\text{gov}}$ to allocative waste and more to transfer. Numbers from Tables~\ref{tab:v3_welfare_ci} and~\ref{tab:v3_strict_invariance}.}
\label{fig:v3_welfare_forest}
\end{figure}

\subsection{Sensitivity to \texorpdfstring{$\lambda$}{lambda}}

For $\lambda \in \{0.20, \lambdaMain, 0.40\}$, the welfare loss runs \welfLossPctLtwentyNp/\welfLossPctLthirtyNp/\welfLossPctLfortyNp{}~percent in non-pharmaceuticals and \welfLossPctLtwentyPh/\welfLossPctLthirtyPh/\welfLossPctLfortyPh{}~percent in pharmaceuticals (Table~\ref{tab:v3_welfare}); each \lambdaInc{} increment in $\lambda$ adds \welfLossDeltaNpPp{}/\welfLossDeltaPhPp{}~pp respectively. Across $\lambda \in [\welfLambdaRangeLo, \welfLambdaRangeHi]$ (Table~\ref{tab:welfare_ranking_lambda}), non-pharma preserves V3 $>$ V0 under both specifications and pharma reverses only under strict invariance. The conditional ranking is therefore driven by the modeling treatment of the post-policy SME pool, not by the choice of $\lambda$.

\paragraph{Bridge to the DiD and strict-invariance reading.} Section~\ref{sec:bridge} attributes the structural-vs-DiD gap to four mechanical sources; eq.~\eqref{eq:welfare} adds a fifth the DiD does not price---the MCPF distortion of \mcpfContribNpPp{}/\mcpfContribPhPp~pp of $p_{S_1}$ in non-pharma/pharma at $\lambda = \lambdaMain$. Strict invariance preserves the direction of the welfare result (\siWelfLossPctNp{}/\siWelfLossPctPh~percent in non-pharma/pharma; Table~\ref{tab:v3_strict_invariance}); the pharma decline relative to the main specification is the footprint of the empirical pre-vs-post difference in $\hat F_c^{\text{SME}}$, an empirical object the framework reports as a bound but does not attribute to a tested selection mechanism.

\subsection{Scaling to vignette and annual aggregate}\label{sec:annual_aggregate}

The pair of furosemida 40\,mg purchases of Section~\ref{sec:data} fixes the magnitudes. With the mean reference price for item~\furosItemId{} at R\$\,\furosUnitPriceBRL{} per unit and a 50{,}000-unit purchase, applying the pharmaceutical-class point estimates yields a total welfare loss on this single purchase of R\$\,\furosTotalLossBRL{} (US\$\,\furosTotalLossUSD)---R\$\,\furosDwlBRL{} of allocative waste, R\$\,\furosTransferBRL{} of implicit transfer to the SME winner, R\$\,\furosMcpfBRL{} of MCPF distortion at $\lambda = \lambdaMain$, summing to \welfLossPctLthirtyPh~percent of $p_{S_1}$ (Table~\ref{tab:v3_welfare}, pharma row).\footnote{R\$/US\$ at R\$\,\fxRate{} (end-of-2018 rate, held fixed across all monetary conversions; sensitivity is mechanical and one-for-one).}

The per-auction number scales. Group-65 procurement in the \structuralWindowM-month Post-policy window totals R\$\,\welfRefOutlayPharma~million in pharmaceutical reference outlays and R\$\,\welfRefOutlayNp~million in non-pharmaceutical reference outlays (own tabulation from the structural sample), or R\$\,\welfRefOutlayPharmaYr~million and R\$\,\welfRefOutlayNpYr~million per year at the Post-policy run-rate. Combining the structural welfare loss per affected auction with the empirical Group-65 spend gives a back-of-envelope annual welfare cost. The headline is the realistic case at the empirical \adherenceGsixtyfivePost~percent state-level adherence rate documented in Section~\ref{sec:institutional}; an upper bound under universal Group-65 coverage under the SME-only rule (\adherenceGsixtyfivePost{}\,$\to$\,100~percent) is reported for context.

Table~\ref{tab:welfare_annual} presents the calculation. The realistic case scales the per-auction welfare share (\welfLossPctLthirtyPh{}~percent of $p_{S_1}$ in pharma, \welfLossPctLthirtyNp{}~percent in non-pharma at $\lambda = \lambdaMain$, from Table~\ref{tab:v3_welfare}) by the \adherenceGsixtyfivePost~percent adherence rate, delivering R\$\,\welfAnnualBRLLo~million (US\$\,\welfAnnualUSDLo~million) per year of welfare loss in Group-65 alone. The upper bound assumes that every Group-65 reference real falls under the SME-only set-aside (counterfactual to the observed regime), giving R\$\,\welfAnnualBRLHi~million (US\$\,\welfAnnualUSDHi~million) per year. Pharma carries the larger share in both cases (R\$\,\welfPharmaAdhBRL~million vs.\ R\$\,\welfNpAdhBRL~million in non-pharma at adherence-adjusted; R\$\,\welfPharmaUbBRL~million vs.\ R\$\,\welfNpUbBRL~million at the upper bound), reflecting the higher per-auction cost share in thin standardized markets.

\begin{table}[!htbp]
\centering
\begin{threeparttable}
\caption{Annual welfare cost of the SME-only set-aside, Group-65 procurement.}
\label{tab:welfare_annual}
\small
\setlength{\tabcolsep}{6pt}
\begin{tabular}{lcc}
\toprule
                            & Adherence-adjusted                    & Upper bound          \\
                            & (\adherenceGsixtyfivePost\% SME-only, observed)   & (full G-65 coverage)  \\
\midrule
Pharma          (R\$~M/yr)  & \welfPharmaAdhBRL   & \welfPharmaUbBRL    \\
Non-pharma      (R\$~M/yr)  & \welfNpAdhBRL       & \welfNpUbBRL        \\
\addlinespace
\textbf{Total Group-65}     & \textbf{R\$\,\welfAnnualBRLLo~M/yr}  & \textbf{R\$\,\welfAnnualBRLHi~M/yr} \\
                            & (US\$\,\welfAnnualUSDLo~M/yr)        & (US\$\,\welfAnnualUSDHi~M/yr)       \\
\bottomrule
\end{tabular}
\begin{tablenotes}\footnotesize
\item Welfare cost per affected auction at $\lambda = \lambdaMain$ is \welfLossPctLthirtyPh~percent of $p_{S_1}$ in pharma and \welfLossPctLthirtyNp~percent in non-pharma (Table~\ref{tab:v3_welfare}). Annual reference outlays from own tabulation of the \structuralWindowM-month Post-policy structural sample (m=\dataOcCutoff--\windowRightStata), annualized to 12-month equivalents: pharma R\$\,\welfRefOutlayPharmaYr~M/yr, non-pharma R\$\,\welfRefOutlayNpYr~M/yr. The adherence-adjusted column (the realistic case and headline of the paper) applies the \adherenceGsixtyfivePost\% state-level Group-65 SME-only adherence rate cited in Section~\ref{sec:institutional}; the upper bound assumes 100\% set-aside coverage (counterfactual to the observed regime) and is reported for context. R\$/US\$ at R\$\,\fxRate.
\end{tablenotes}
\end{threeparttable}
\end{table}

\paragraph{Sensitivity to the adherence rate.}\label{sec:adherence_sensitivity} The structural content of the welfare exercise is the per-auction estimate ($\welfTotalLossLthirtyPh \times p^{\mathrm{ref}}$ in pharma, $\welfTotalLossLthirtyNp \times p^{\mathrm{ref}}$ in non-pharma at $\lambda = \lambdaMain$; Table~\ref{tab:v3_welfare}), identified by the model and invariant to institutional scaling. The adherence rate translates that object to an annual aggregate linearly: across the realistic range $[\welfAdherenceLoPct\%, \welfAdherenceHiPct\%]$ the annual welfare cost spans R\$\,\welfRangeLoBRL--\welfRangeHiBRL~million (US\$\,\welfRangeLoUSD--\welfRangeHiUSD~million), with the empirical \adherenceGsixtyfivePost~percent rate delivering the R\$\,\welfAnnualBRLLo~million central case. The \adherenceGsixtyfivePost~percent rate pools class-6531 pharmaceuticals (lower-adhering because the three-eligible-SMEs requirement binds more often in oligopolistic markets) with non-pharma supplies; the adherence-adjusted headline is a conservative reading on Group-65 alone, with the platform's other SME-eligible items and the federal extension under LC \leiNova{} scaling proportionally above.

\begin{table}[!htbp]
\centering
\begin{threeparttable}
\caption{Annual welfare cost: sensitivity to SME-only adherence rate within Group 65.}
\label{tab:welfare_adherence_sensitivity}
\small
\setlength{\tabcolsep}{4pt}
\begin{tabular}{lcccccc}
\toprule
Adherence rate & 30\% & \textbf{43\%} & 55\% & 70\% & 85\% & 100\% \\
\midrule
Pharma           (R\$~M/yr) & 22 & \textbf{31} & 40 & 51 & 62 & 73 \\
Non-pharma       (R\$~M/yr) & 16 & \textbf{24} & 30 & 38 & 47 & 55 \\
\addlinespace
\textbf{Total Group-65} & \textbf{R\$\,38} & \textbf{R\$\,55} & \textbf{R\$\,70} & \textbf{R\$\,89} & \textbf{R\$\,109} & \textbf{R\$\,128} \\
 & (US\$\,11) & \textbf{(US\$\,16)} & (US\$\,20) & (US\$\,25) & (US\$\,31) & (US\$\,37) \\
\bottomrule
\end{tabular}
\begin{tablenotes}\footnotesize
\item Welfare cost per affected auction held fixed at the $\lambda = 0.30$ point estimates of Table~\ref{tab:v3_welfare}: $0.308 \times p^{\mathrm{ref}}$ in pharma, $0.204 \times p^{\mathrm{ref}}$ in non-pharma. Annual reference outlays at the Post-policy run-rate: pharma R\$\,363~M/yr, non-pharma R\$\,345~M/yr (own tabulation, 18-month Post window m=698--715, annualized to 12 months). Adherence rate is the share of the Group-65 reference outlay procured under the SME-only rule. Baseline column (43\%) shown in bold; this rate is the empirical Post-period state-level Group-65 adherence cited in Section~\ref{sec:institutional}. R\$/US\$ at R\$\,3.50 (end-2018).
\end{tablenotes}
\end{threeparttable}
\end{table}

